% SPDT paper -- IEEE conference manuscript
% Official IEEE template: IEEEtran.cls v1.8+ (conference mode)
% Compile: latexmk -pdf main.tex  (see paper/README.md)
\documentclass[conference]{IEEEtran}
\IEEEoverridecommandlockouts

\usepackage{cite}                 % IEEE-style numeric citation compression
\usepackage{amsmath,amsthm,amsfonts}
\usepackage{booktabs}
\usepackage{array}
\usepackage{textcomp}
\usepackage{xurl}                 % line-breakable URLs in references
\usepackage{graphicx}

% ---------------------------------------------------------------------------
% TODO(AUTHORSHIP) -- SIX FACULTY CO-AUTHORS REQUIRED BY THE COURSE.
% Identities have NOT been confirmed by the user. Do not fabricate names.
% After confirmation, extend the \author block with:
%   \and
%   \IEEEauthorblockN{<Full Name>}
%   \IEEEauthorblockA{<Department>\\<University>\\City, Country\\email}
% The first author (student) must remain first.
% ---------------------------------------------------------------------------

\title{Stability-Penalized Split Selection for Decision Tree Classification:\\
A Reproducible Multi-Dataset Evaluation}

\author{\IEEEauthorblockN{Rofiq Fauzi}
\IEEEauthorblockA{Magister Artificial Intelligence\\
Universitas Dian Nuswantoro\\
Semarang, Indonesia\\
p36202600047@mhs.dinus.ac.id}}

\begin{document}

\maketitle

\begin{abstract}
The CART decision tree is fast and interpretable, but its structure is
unstable: small changes in the training data can alter the selected split
and, ultimately, the predictions. This study evaluates a simple candidate
formulation, named Stability-Penalized Decision Tree (SPDT), in which each
candidate split is scored as Gini gain minus a weighted penalty on the
variance of that gain over bootstrap resamples of the node data.
A literature search up to 22 October 2026 found no formula-level
equivalent of this criterion, but the tree-stability research space is
crowded; the formulation is therefore positioned as a candidate, not as a
novelty claim. The evaluation uses eight UCI tabular classification
datasets with repeated nested cross-validation (5 folds $\times$ 3 repeats
on three seeds), hyperparameter tuning in the inner loop, and Macro-F1 as
the primary metric. Results: SPDT versus custom CART ($\lambda{=}0$)
records six wins, zero ties, and two losses with a mean Macro-F1
difference of $+0.0058$. The Friedman test gives
$\chi^2(2){=}4.0000$, $p{=}0.1353$, Kendall's $W{=}0.2500$ --- insufficient
evidence to reject the null hypothesis, so confirmatory post-hoc
comparisons were not run. Ablations show that the penalty strength
$\lambda$ and the number of resamples $B$ barely affect Macro-F1, tree
depth constraints matter far more, and the prediction-stability measure of
SPDT is slightly worse than that of custom CART at roughly $40\times$ the
fitting time. These findings are an honest null result: across the eight
small-to-medium datasets tested, the bootstrap gain-variance penalty
provides no benefit sufficient to outweigh the simplicity of CART, without
ruling out benefits in other data regimes.
\end{abstract}

\begin{IEEEkeywords}
Decision trees, split criterion, bootstrap stability, gain variance,
nested cross-validation, nonparametric statistical comparison.
\end{IEEEkeywords}

\section{Introduction}
\IEEEPARstart{D}{ecision} trees are nonparametric classifiers that
recursively partition the feature space. The CART split criterion selects
the candidate (feature, threshold) that maximizes the reduction of Gini
impurity \cite{breiman1984cart,quinlan1986induction}. However, tree
structure is known to be unstable: small perturbations of the training
data can change the selected split \cite{li2002instability}, and
consequently the predictions --- the very phenomenon behind bagging and
random forests \cite{breiman1996bagging,breiman2001randomforests,dietterich2000ensemble}.

Various approaches have reduced this instability at different levels:
significance filtering, structural restriction, ensemble aggregation, or
stable model selection \cite{turney1995stability,briand2009stability,lee2024anls,mirzamomen2017stable}.
What is rarely tested directly is the narrower question:
\textbf{can a bootstrap variance penalty on the split gain, subtracted
directly from the split score, improve out-of-sample Macro-F1 and/or
prediction stability compared with standard CART under repeated
resampling, at acceptable runtime cost?}

Null hypothesis: no systematic improvement over the baseline across
independent datasets. No hypothesis is made that SPDT must win. The
claimed contributions are limited to (i)~an explicit candidate
formulation, (ii)~a reproducible implementation and audit, (iii)~a
leakage-free nested-CV evaluation, and (iv)~honest ablation and
statistical analysis of both positive and negative outcomes.

\section{Related Work}
\subsection{Decision Tree Stability}
Structural tree instability was formalized by Li and Clarke
\cite{li2002instability}. Quantitative stability and tree-similarity
measures were developed by Turney \cite{turney1995stability} and Briand
et al.\ \cite{briand2009stability}, who also evaluated bootstrap-based
node-level stabilization (the NLS/Dannegger lineage). Lee et al.\
\cite{lee2024anls} proposed ANLS, which runs bootstrap iterations at each
node and selects the most frequently chosen predictor --- the closest
resampling-per-node mechanism to this study, but its aggregation is
frequency voting rather than a gain-variance penalty. Mirzamomen and
Kangavari \cite{mirzamomen2017stable} stabilize trees through
multi-attribute split tests.

\subsection{Alternative Split Criteria}
Two-term Tsallis entropy criteria \cite{wang2017tsallis}, label-noise-robust
criteria \cite{yang2019pairgain}, and bias-corrected Gini split selection
\cite{strobl2007bias} show that the ``modify the split score'' space is
active. Hoeffding-bound data-stream trees \cite{domingos2000vfdt} decide
splits only when the gain estimate is separated from the runner-up by a
confidence interval, and gradient-boosted trees regularize splitting with
a minimum-split-gain threshold \cite{chen2016xgboost}. Delise
\cite{delise2026era} scores split consistency across environments in
boosting. Penalizing a resampled score by its variance is also known
outside trees (the mean--variance pattern in model selection). None of the
works we verified contains the formula
$\mathrm{Gain}-\lambda\,\mathrm{Var}_{\text{bootstrap}}[\mathrm{Gain}]$;
nevertheless, a negative search result is not proof of novelty, and the
formulation remains a candidate.

\section{Method}
\subsection{Proposed Split Criterion}
For a node with dataset $D$ (training portion only), class
probabilities $p(c\mid D)$, and $C$ classes, the Gini impurity is
\begin{equation}
\operatorname{Gini}(D) = 1 - \sum_{c=1}^{C} p(c\mid D)^2 .
\label{eq:gini}
\end{equation}
For a candidate split $s$ partitioning $D$ into $D_L$ and $D_R$, the CART
gain is
\begin{equation}
\operatorname{Gain}(s,D) = \operatorname{Gini}(D)
- \tfrac{|D_L|}{|D|}\operatorname{Gini}(D_L)
- \tfrac{|D_R|}{|D|}\operatorname{Gini}(D_R).
\label{eq:gain}
\end{equation}
From $D$ we draw $B$ bootstrap resamples $D^{(b)}$ (sampling with
replacement, node rows only). With the candidate $(s)$ fixed, the
resampled gain variance is
\begin{equation}
V(s) = \operatorname{Var}_{b}\!\left[\operatorname{Gain}(s,D^{(b)})\right],
\qquad b = 1,\dots,B,
\label{eq:var}
\end{equation}
and the proposed score is
\begin{equation}
\boxed{\;\operatorname{Score}(s) = \operatorname{Gain}(s,D) - \lambda\, V(s)\;}
\label{eq:score}
\end{equation}
where $\lambda \geq 0$ controls the penalty strength. Here $D$ is the
current node's training data, $s$ is a (feature, threshold) candidate,
$B$ is the number of bootstrap resamples, and $\lambda$ weights the
penalty; $\lambda{=}0$ disables the penalty entirely. Candidate thresholds
are midpoints between sorted distinct adjacent values;
$\mathrm{min\_samples\_leaf}$ and $\mathrm{min\_samples\_split}$ are applied
to original node counts; an empty child partition inside a resample
contributes zero; leaves predict the majority class with smallest-label
tie-breaking; bootstrap seeds derive deterministically from
$(\mathrm{random\_state}, \mathrm{node\ id})$. For $\lambda{=}0$ the
estimator degenerates to a custom CART reference, verified identical to an
independent reference implementation in unit tests. A separate exploratory
variant normalizing $V(s)$ by $\operatorname{Gain}(s,D)^2+\varepsilon$ is
documented but is not the main method. An important limitation is that
$V(s)$ is scale-dependent on the gain.

\subsection{Evaluation Design}
We use eight UCI tabular classification datasets (Table~\ref{tab:inventory}).
The protocol is repeated stratified nested CV: an outer loop of 5 folds
$\times$ 3 repeats for each of the seeds $[42,123,2026]$ (45 evaluations
per dataset--method, identical splits for all methods), and an inner
stratified 3-fold loop for tuning with deterministically derived seeds.
Preprocessing is median imputation fitted on the training portion only.
The primary metric is Macro-F1; secondary measures are balanced accuracy,
accuracy, fit/predict time, node count, and tree depth. Prediction
stability is measured separately as mean pairwise disagreement among
models trained on ten 90\% subsamples of each outer training fold,
evaluated on the same outer test fold, and is reported separately from
accuracy. The primary comparison set is sklearn CART, custom CART
($\lambda{=}0$), and SPDT; a random forest serves as a non-equivalent
contextual benchmark outside the omnibus test.

\subsection{Statistical Analysis}
The unit of inference is the dataset (averaged over the 45 outer folds),
never the individual fold. The three primary methods are compared with the
Friedman test using datasets as blocks; only when the omnibus test
provides grounds for it are paired Wilcoxon signed-rank tests with Holm
correction and rank-biserial effect sizes run. $p \geq 0.05$ is never
interpreted as evidence of equivalence.

\section{Results and Discussion}
Table~\ref{tab:performance} reports per-dataset Macro-F1 (mean $\pm$ SD
over the 45 outer folds, descriptive only). Table~\ref{tab:ranks} reports
mean ranks, and Table~\ref{tab:friedman} reports the omnibus test. SPDT
versus custom CART records six wins, zero ties, and two losses with mean
Macro-F1 difference $+0.0058$. The Friedman test gives
$\chi^2(2)=4.0000$, $p=0.1353$, Kendall's $W=0.2500$ --- insufficient
evidence to reject the null hypothesis, so confirmatory post-hoc
comparisons were not run. The exploratory SPDT--custom CART Wilcoxon
comparison gives a raw $p=0.078$, which does not survive the 5\% level
and must not be read as significant; with only eight datasets, statistical
power is limited. These results are reported as they stand.

% AUTO-GENERATED from results/ by make_ieee_tables.py
% Do not hand-edit numbers; regenerate instead.

\begin{table}[!t]
\caption{Dataset Inventory}
\label{tab:inventory}
\centering
\begin{tabular}{@{}lrrrl@{}}
\toprule
Dataset & $N$ & Features & Classes & Source \\
\midrule
banknote & 1372 & 4 & 2 & ucimlrepo \\
glass & 214 & 9 & 6 & ucimlrepo \\
ionosphere & 351 & 34 & 2 & ucimlrepo \\
iris & 150 & 4 & 3 & ucimlrepo \\
seeds & 210 & 7 & 3 & UCI zip \\
sonar & 208 & 60 & 2 & ucimlrepo \\
wdbc & 569 & 30 & 2 & ucimlrepo \\
wine & 178 & 13 & 3 & ucimlrepo \\
\bottomrule
\end{tabular}
\end{table}

\begin{table*}[!t]
\caption{Macro-F1 Mean $\pm$ SD (Descriptive Fold-Level Scores) per Dataset and Method}
\label{tab:performance}
\centering
\begin{tabular*}{\textwidth}{@{\extracolsep{\fill}}lccccc@{}}
\toprule
Dataset & CART-skl & CART-cust & SPDT & RF & SPDT$-$CART \\
\midrule
banknote & 0.982 $\pm$ 0.008 & 0.981 $\pm$ 0.007 & 0.981 $\pm$ 0.007 & 0.994 $\pm$ 0.004 & +0.0004 \\
glass & 0.620 $\pm$ 0.097 & 0.609 $\pm$ 0.093 & 0.626 $\pm$ 0.093 & 0.742 $\pm$ 0.075 & +0.0166 \\
ionosphere & 0.873 $\pm$ 0.034 & 0.881 $\pm$ 0.038 & 0.879 $\pm$ 0.038 & 0.929 $\pm$ 0.037 & -0.0027 \\
iris & 0.949 $\pm$ 0.033 & 0.945 $\pm$ 0.038 & 0.944 $\pm$ 0.037 & 0.955 $\pm$ 0.034 & -0.0008 \\
seeds & 0.897 $\pm$ 0.032 & 0.906 $\pm$ 0.036 & 0.913 $\pm$ 0.034 & 0.923 $\pm$ 0.038 & +0.0066 \\
sonar & 0.719 $\pm$ 0.065 & 0.706 $\pm$ 0.073 & 0.725 $\pm$ 0.061 & 0.833 $\pm$ 0.065 & +0.0189 \\
wdbc & 0.925 $\pm$ 0.031 & 0.925 $\pm$ 0.030 & 0.928 $\pm$ 0.030 & 0.956 $\pm$ 0.023 & +0.0037 \\
wine & 0.905 $\pm$ 0.047 & 0.905 $\pm$ 0.044 & 0.909 $\pm$ 0.046 & 0.981 $\pm$ 0.024 & +0.0040 \\
\bottomrule
\end{tabular*}
\end{table*}

\begin{table}[!t]
\caption{Mean Macro-F1 Rank Across Datasets}
\label{tab:ranks}
\centering
\begin{tabular}{@{}lr@{}}
\toprule
Method & Mean rank \\
\midrule
SPDT & 1.500 \\
CART-skl & 2.000 \\
CART-cust & 2.500 \\
\bottomrule
\end{tabular}
\end{table}

\begin{table}[!t]
\caption{Dataset-Level Statistical Test (Friedman, $D{=}8$ Datasets, 3 Methods)}
\label{tab:friedman}
\centering
\begin{tabular}{@{}lr@{}}
\toprule
Statistic & Value \\
\midrule
$\chi^2$ & 4.0000 \\
Degrees of freedom & 2 \\
$p$ & 0.1353 \\
Kendall's $W$ & 0.2500 \\
Decision & post-hoc not run \\
\bottomrule
\end{tabular}
\end{table}

\begin{table}[!t]
\caption{Computational Cost (per Outer Fold, Averaged Over Datasets)}
\label{tab:runtime}
\centering
\begin{tabular}{@{}lrrrrr@{}}
\toprule
Method & Fit (s) & Fit median (s) & Predict (s) & Nodes & Depth \\
\midrule
CART-cust & 0.5340 & 0.4163 & 0.0009 & 26.6722 & 5.1000 \\
CART-skl & 0.2693 & 0.2605 & 0.0010 & 27.4222 & 5.2917 \\
RF & 1.1169 & 1.0863 & 0.0573 & - & - \\
SPDT & 21.4302 & 14.6491 & 0.0010 & 28.5944 & 5.4611 \\
\bottomrule
\end{tabular}
\end{table}

\begin{table}[!t]
\caption{Prediction Stability (Mean Pairwise Disagreement; Not an Accuracy Metric)}
\label{tab:stability}
\centering
\begin{tabular}{@{}lrr@{}}
\toprule
Method & Mean disagreement & SD \\
\midrule
CART-cust & 0.0928 & 0.0853 \\
CART-skl & 0.0983 & 0.0885 \\
RF & 0.0271 & 0.0337 \\
SPDT & 0.1014 & 0.0934 \\
\bottomrule
\end{tabular}
\end{table}

\begin{table}[!t]
\caption{Ablation on Identical Outer Splits}
\label{tab:ablation}
\centering
\begin{tabular}{@{}lrrrr@{}}
\toprule
Variant & Macro F1 & SD & Fit (s) & Nodes \\
\midrule
A0 ($\lambda$=0) & 0.8599 & 0.1270 & 0.0201 & 37.6611 \\
A1 ($\lambda$=.1) & 0.8598 & 0.1256 & 0.0860 & 37.6111 \\
A1 ($\lambda$=1) & 0.8601 & 0.1255 & 0.0862 & 37.5278 \\
A2 ($B$=20) & 0.8601 & 0.1251 & 0.1136 & 37.4556 \\
A2 ($B$=5) & 0.8613 & 0.1243 & 0.0711 & 37.4667 \\
A3 d3 l1 & 0.8357 & 0.1592 & 0.0371 & 12.8722 \\
A3 d3 l5 & 0.8353 & 0.1546 & 0.0373 & 11.8611 \\
A3 d8 l1 & 0.8592 & 0.1272 & 0.0821 & 35.2500 \\
A3 d8 l5 & 0.8521 & 0.1381 & 0.0820 & 24.7389 \\
\bottomrule
\end{tabular}
\end{table}

\begin{table}[!t]
\caption{SPDT Hyperparameter Selection Frequency (45 Outer Folds per Dataset)}
\label{tab:hyper}
\centering
\begin{tabular}{@{}llrr@{}}
\toprule
Parameter & Value & Frequency & Share \\
\midrule
max\_depth & 5.0 & 140 & 0.3889 \\
max\_depth & 8.0 & 102 & 0.2833 \\
max\_depth & 3.0 & 92 & 0.2556 \\
min\_samples\_leaf & 1.0 & 201 & 0.5583 \\
min\_samples\_leaf & 2.0 & 97 & 0.2694 \\
min\_samples\_leaf & 5.0 & 62 & 0.1722 \\
min\_samples\_split & 2.0 & 252 & 0.7000 \\
min\_samples\_split & 5.0 & 69 & 0.1917 \\
min\_samples\_split & 10.0 & 39 & 0.1083 \\
n\_bootstrap & 5.0 & 191 & 0.5306 \\
n\_bootstrap & 20.0 & 87 & 0.2417 \\
n\_bootstrap & 10.0 & 82 & 0.2278 \\
penalty\_lambda & 10.0 & 233 & 0.6472 \\
penalty\_lambda & 0.0 & 61 & 0.1694 \\
penalty\_lambda & 0.01 & 41 & 0.1139 \\
penalty\_lambda & 1.0 & 22 & 0.0611 \\
penalty\_lambda & 0.1 & 3 & 0.0083 \\
\bottomrule
\end{tabular}
\end{table}


A second finding concerns stability. The bootstrap-variance penalty does
\emph{not} reduce prediction disagreement: SPDT is slightly worse
($0.101$) than custom CART ($0.093$) and sklearn CART ($0.098$), while
the random forest remains the most stable ($0.027$) through ensemble
aggregation (Table~\ref{tab:stability}). The motivating promise of the
method --- better prediction stability --- therefore does not materialize
on this benchmark; a gain-variance penalty is no substitute for
ensembling.

Third, the identical-split ablations (Table~\ref{tab:ablation}) show that
neither the penalty strength $\lambda$ nor the number of resamples $B$
materially changes Macro-F1 (all $\lambda$ and $B$ variants fall in the
$0.859$--$0.861$ range), whereas depth/leaf constraints matter far more
(depth 3 drops to about $0.836$). Inner tuning frequently selects a large
$\lambda$ (10 on 64.7\% of folds) and a small $B$ (5 on 53.1\%), per
Table~\ref{tab:hyper}, indicating the penalty is often selected but has a
small effect on accuracy. The cost is real: SPDT fitting time is roughly
$40\times$ that of custom CART (Table~\ref{tab:runtime}). Overall, the
effectiveness of the bootstrap gain-variance penalty is dataset-dependent
and small in magnitude; simple impurity-based splitting remains
preferable for these small-to-medium datasets.

\section{Conclusion}
This study evaluated the SPDT candidate formulation reproducibly, with
leakage-free repeated nested CV, controlled ablation on identical splits,
and dataset-level statistical analysis on eight UCI datasets.
Descriptively, SPDT achieved the best mean rank and won six of eight
paired comparisons against custom CART, but the Friedman test provided
insufficient evidence of a difference ($\chi^2(2)=4.0000$, $p=0.1353$),
so there is no confirmatory basis for claiming improvement. Moreover, the
prediction-stability measure did not improve and was slightly worse,
while compute cost rose roughly $40\times$. Under the tested conditions,
the simplicity of CART remains preferable. Limitations: eight small
datasets, Gini impurity only, a single disagreement-based stability proxy,
a formula-level novelty claim that is not established, and the absence of
the closest literature criteria (ANLS, minimax split) as comparators. The
most valuable future work is direct comparison with node-level bootstrap
stabilization from the literature, a theoretically grounded penalty
normalization by gain scale, and evaluation on larger benchmarks.

\bibliographystyle{IEEEtran}
\bibliography{references}

\end{document}
